% ============ 31 Lefschetz 铅笔的几何 (原书 pp.184-185) ============
\chapter{Lefschetz 铅笔的几何}\label{sec:31}
\index{Lefschetz 铅笔}
\index{对偶射影空间}
\index{铅笔}
\index{普通二重点}
\index{对偶簇}
\index{爆破}

本节中，我们研究如何把簇纤维化到 $\mathbb{P}^{1}$ 上，使纤维只有非常简单的奇异性；下一节我们用这个纤维化研究簇的上同调.这种研究簇上同调的方法在复情形可以追溯到 Lefschetz.本节中，我们在代数闭域 $k$ 上工作.

\subsection*{定义}

线性形式 $H = \sum_{i=0}^{m} a_i T_i$ 定义 $\mathbb{P}^{m}$ 中的超平面，两个线性形式定义同一个超平面，当且仅当一个是另一个的非零倍.于是 $\mathbb{P}^{m}$ 中的超平面构成射影空间，称为对偶射影空间 $\check{\mathbb{P}}^{m}$.

$\check{\mathbb{P}}^{m}$ 中的直线 $D$ 称为 $\mathbb{P}^{m}$ 中的超平面铅笔.若 $H_0$ 与 $H_1$ 是 $D$ 中任意两个互异超平面，则铅笔由所有形如 $\alpha H_0 + \beta H_1$ 的超平面组成($(\alpha \colon \beta) \in \mathbb{P}^{1}(k)$).若 $P \in H_0 \cap H_1$，则 $P$ 位于铅笔的每个超平面中---铅笔的轴 $A$ 定义为这样的 $P$ 的集合.于是
\[
A = H_0 \cap H_1 = \bigcap_{t \in D} H_t.
\]
铅笔的轴是 $\mathbb{P}^{m}$ 中余维数 2 的线性子簇，而铅笔中的超平面恰好是含轴的超平面.过 $\mathbb{P}^{m}$ 中不在 $A$ 上的任何点，恰好有一个铅笔中的超平面.于是，应该想象铅笔中的超平面绕轴旋转扫出 $\mathbb{P}^{m}$.

设 $X$ 是维数 $d \geq 2$ 的非奇异射影簇，把它嵌入某个射影空间 $\mathbb{P}^{m}$.所谓嵌入的平方，我们指 $X \hookrightarrow \mathbb{P}^{m}$ 与 Veronese 映射(AG 6.20)
\[
(x_0 \colon \ldots \colon x_m) \mapsto (x^{2}_0 \colon \ldots \colon x_i x_j \colon \ldots \colon x^{2}_m) \colon \mathbb{P}^{m} \to \mathbb{P}^{\frac{(m+2)(m+1)}{2}}
\]
的复合.

\begin{definition}\label{def:31.1}
$\check{\mathbb{P}}^{m}$ 中的直线 $D$ 称为 $X \subset \mathbb{P}^{m}$ 的 Lefschetz 铅笔，如果
\begin{enumerate}
  \item 铅笔 $(H_t)_{t \in D}$ 的轴 $A$ 横截相交 $X$；
  \item 对 $D$ 的某个开稠密子集 $U$ 中的所有 $t$，$X$ 的超平面截 $X_t \overset{\mathrm{def}}{=} X \cap H_t$ 非奇异；
  \item 对 $t \notin U$，$X_t$ 只有单个奇点，且该奇点是普通二重点.
\end{enumerate}
条件 (a) 意味着对任何闭点 $P \in A \cap X$，$\operatorname{Tgt}_{P}(A) \cap \operatorname{Tgt}_{P}(X)$ 在 $\operatorname{Tgt}_{P}(X)$ 中有余维数 2.
\end{definition}

(b) 中的交 $X \cap H_t$ 应在概形意义下取，即若 $X$ 由齐次理想 $\mathfrak{a}$ 定义，则 $X \cap H_t$ 由 $\mathfrak{a} + (H_t)$ 定义.条件 (b) 意味着 $X_t$ 既约且作为代数簇非奇异.

维数 $d$ 的簇 $X$ 上的点 $P$ 称为普通二重点，如果 $P$ 处的切锥同构于 $\mathbb{A}^{d+1}$ 中由非退化二次型 $Q(T_1, \ldots, T_{d+1})$ 定义的子簇，或者等价地，如果
\[
\widehat{\mathcal{O}}_{X, P} \cong k[[T_1, \ldots, T_{d+1}]]/(Q(T_1, \ldots, T_{d+1})).
\]

\begin{theorem}\label{thm:31.2}
对 $X$ 存在 Lefschetz 铅笔(在可能把 $X$ 的射影嵌入换成其平方后).
\end{theorem}

\begin{proof}
(梗概) 设 $Y \subset X \times \check{\mathbb{P}}^{m}$ 是点对 $(x, H)$ 组成的闭簇，其中 $H$ 含 $X$ 在 $x$ 处的切空间.例如，若 $X$ 在 $\mathbb{P}^{m}$ 中有余维数 1，则 $(x, H) \in Y$，当且仅当 $H$ 是 $x$ 处的切空间.一般地，
\[
(x, H) \in Y \Longleftrightarrow x \in H \text{ 且 } H \text{ 在 } x \text{ 处不横截相交 } X.
\]
$Y$ 在投影映射 $X \times \check{\mathbb{P}}^{m} \to \check{\mathbb{P}}^{m}$ 下的像称为 $X$ 的对偶簇 $\check{X}$.

$Y \to X$ 在 $x$ 上的纤维由含 $x$ 处切空间的超平面组成，这些超平面构成 $\check{\mathbb{P}}^{m}$ 中维数 $m - (\dim X + 1)$ 的不可约子簇；由此推出 $Y$ 不可约，完备，维数为 $m - 1$(见 AG 10.11)，且 $\check{X}$ 不可约，完备，在 $\check{\mathbb{P}}^{m}$ 中余维数 $\geq 1$(除非 $X = \mathbb{P}^{m}$，这时它为空).映射 $\varphi \colon Y \to \check{X}$ 在 $(x, H)$ 处非分歧，当且仅当 $x$ 是 $X \cap H$ 的普通二重点(见 SGA 7, XVII 3.7).要么 $\varphi$ 一般非分歧，要么把嵌入换成其平方(于是不是用超平面，而是用二次超曲面)后变成这样(同上，3.7).我们可以假定如此，然后(同上，3.5)可以证明：对 $H \in \check{X} \setminus \check{X}_{\mathrm{sing}}$，$X \cap H$ 只有单个奇点，且该奇点是普通二重点.这里 $\check{X}_{\mathrm{sing}}$ 是 $\check{X}$ 的奇异轨迹.

由 Bertini 定理(Hartshorne II.8.18)存在超平面 $H_0$，使得 $H_0 \cap X$ 不可约且非奇异.由于过 $H_0$ 的直线构成 $m - 1$ 维空间，且至多有 $m - 2$ 维的族与 $\check{X}_{\mathrm{sing}}$ 相交，可以选取 $H_1$，使得联结 $H_0$ 与 $H_1$ 的直线 $D$ 不与 $\check{X}_{\mathrm{sing}}$ 相交.则 $D$ 是 $X$ 的 Lefschetz 铅笔.
\end{proof}

\begin{theorem}\label{thm:31.3}
设 $D = (H_t)$ 是轴为 $A = \bigcap H_t$ 的 $X$ 的 Lefschetz 铅笔.则存在簇 $X^{*}$ 与映射
\[
X \leftarrow X^{*} \xrightarrow{\varphi} D,
\]
使得：
\begin{enumerate}
  \item 映射 $X^{*} \to X$ 是 $X$ 沿 $A \cap X$ 的爆破；
  \item $X^{*} \to D$ 在 $t$ 上的纤维是 $X_t = X \cap H_t$.
\end{enumerate}
此外，$\varphi$ 紧合，平坦，且有截面.
\end{theorem}

\begin{proof}
(梗概) 过 $X \setminus A \cap X$ 的每个点 $x$，恰好有一个 $D$ 中的 $H_x$.映射
\[
\varphi' \colon X \setminus A \cap X \to D, \qquad x \mapsto H_x,
\]
正则.取其图象 $\Gamma_{\varphi'}$ 在 $X \times D$ 中的闭包；它将是 $\varphi$ 的图象.
\end{proof}

\begin{remark}\label{rem:31.4}
奇异的 $X_t$ 可能是可约的.例如，若 $X$ 是 $\mathbb{P}^{3}$ 中的二次曲面，则对所有 $t$，$X_t$ 是 $\mathbb{P}^{2}$ 中次数为 2 的曲线，而这样的曲线奇异，当且仅当它可约(看亏格公式).然而，若把嵌入 $X \hookrightarrow \mathbb{P}^{m}$ 换成其立方，这个问题将永不出现.
\end{remark}

\textbf{参考文献} 我知道的唯一现代参考文献是 SGA 7, Expos\'e XVII.(也许某天我会把它放进 AG.)
